% Kapitel 2 – wird von masterarbeit.tex / abgabe.tex eingebunden
\chapter{Stand der Technik}
\label{ch:sota}
% ==============================================================================

Dieses Kapitel gibt einen Überblick über die Blickschätzung mit normalen
RGB-Kameras. Zuerst erkläre ich die Grundlagen (\cref{sec:grundlagen}) und die
besonderen Probleme von Webcams (\cref{sec:webcam}). Danach stelle ich wichtige
Datensätze (\cref{sec:datensaetze}), Netzarchitekturen
(\cref{sec:architekturen}) und vortrainierte Modelle (\cref{sec:pretrained})
vor. Anschließend geht es um Personalisierung, Domain Shift, Test-Time
Adaptation und Datenqualität. Am Ende leite ich die Forschungslücken ab, die
diese Arbeit bearbeitet.

% ------------------------------------------------------------------------------
\section{Grundlagen der erscheinungsbasierten Blickschätzung}
\label{sec:grundlagen}
% ------------------------------------------------------------------------------
Eye-Tracking bestimmt aus Bildern der Augen oder des Gesichts, wohin eine
Person schaut. Es gibt zwei Arten von Ausgaben: die dreidimensionale
Blickrichtung (meist zwei Winkel, Pitch und Yaw) oder den zweidimensionalen
\emph{Point of Regard} (PoR). Der PoR ist der Punkt auf dem Bildschirm, auf den
die Person schaut. Diese Arbeit schätzt den 2D-PoR aus Bildern einer normalen
Webcam.

\subsection{Klassische Verfahren}
\label{subsec:klassisch}
Erscheinungsbasierte Verfahren lernen direkt eine Abbildung vom Bild zum
Blickziel. Sie brauchen kein geometrisches Modell des Auges. Frühe Verfahren
haben von Hand entworfene Bild-Features verwendet, zum Beispiel Local Binary
Patterns (LBP), Histograms of Oriented Gradients (HoG) oder SIFT. Diese
Features wurden mit einfachen Regressoren kombiniert, etwa mit linearer
Regression, Random Forests oder Support-Vector-Regression (SVR). Solche
Verfahren rechnen schnell. Bei Kopfbewegungen und wechselndem Licht werden sie
aber deutlich ungenauer, weil die Features diese Änderungen nicht gut von der
Blickinformation trennen.

\subsection{Erste Deep-Learning-Verfahren}
\label{subsec:dl_pioniere}
Mit Convolutional Neural Networks (CNNs) kann man Gesichts- und Augenbilder
direkt auf Bildschirmpositionen abbilden. Krafka et al.~\cite{krafka2016itracker}
haben mit \emph{iTracker} eines der ersten CNNs für mobiles Eye-Tracking ohne
zusätzliche Hardware vorgestellt. Das Modell verarbeitet vier Eingaben in
getrennten Zweigen: das linke Auge, das rechte Auge, das Gesicht und ein
binäres \emph{Face Grid} (wo das Gesicht im Bild liegt). Danach werden die
Features zusammengeführt. Ohne Kalibrierung liegt der Fehler bei
\SI{1,71}{\centi\metre} auf Smartphones und bei \SI{2,53}{\centi\metre} auf
Tablets. Für die Personalisierung trainieren die Autoren eine SVR pro Nutzer
auf den Features der Schicht \code{fc1}. Dafür nutzen sie Aufnahmen von bis zu
13 festen Kalibrierpunkten. Mit allen 13 Punkten sinkt der Fehler auf
\SI{1,34}{\centi\metre} bzw.\ \SI{2,12}{\centi\metre}. Mit sehr wenigen Punkten
kommt es bei der SVR aber leicht zu \emph{Overfitting} (dt.\ Überanpassung):
Das Modell lernt die wenigen Trainingsdaten fast auswendig und wird dadurch
auf neuen Bildern schlechter. Das
Modell läuft in Echtzeit auf mobilen Geräten.

Ein wichtiger Beitrag dieser Arbeit ist der Datensatz GazeCapture
(\cref{sec:datensaetze}). Er wurde per \emph{Crowdsourcing} auf Smartphones
und Tablets unter Alltagsbedingungen gesammelt: Viele Freiwillige haben die
Daten selbst mit einer App auf ihrem eigenen Gerät aufgenommen. Er dient
bis heute als Benchmark und als Quelle für das Pretraining.

\paragraph{Bedeutung für diese Arbeit.}
iTracker hat drei Ideen eingeführt, die bis heute wichtig sind: (1) den
Datensatz GazeCapture, (2) die Personalisierung mit Kalibrierdaten und (3) das
Zusammenführen von Features aus mehreren Bildbereichen. Diese Arbeit nutzt
dieselben Daten. Statt AlexNet und SVR verwendet sie aber ein ResNet-18
\cite{he2016resnet} mit einem Regression Head und L1-Loss. Der
Regression Head ist hier ein Multi-Layer Perceptron (MLP), also ein kleines Netz aus
vollständig verbundenen Schichten, das aus den Bild-Features des ResNet-18 die
Koordinaten $(x, y)$ auf dem Bildschirm berechnet. Außerdem
bewertet sie die Basismodelle mit einem personendisjunkten Split
(\code{norm\_subject}) und nicht nur mit einer zufälligen Aufteilung
(\code{norm\_random}). Falsche Labels behandelt sie schon bei der Aufnahme: Das
Aufnahmeprogramm speichert für jedes Bild die Phase des Ziels
(\code{WAIT} oder \code{MOVE}).

\subsection{Einteilung der Verfahren}
\label{subsec:taxonomie}
Man kann Verfahren zur Blickschätzung in drei Klassen einteilen:
\begin{enumerate}
  \item \textbf{Modellbasierte (geometrische) Verfahren} nutzen ein 3D-Modell
        des Augapfels, Pupillenellipsen und Lichtreflexe auf der Hornhaut. Sie
        brauchen meist hochaufgelöste Bilder und Infrarotlicht. Bei normalem
        Tageslicht machen sie oft Fehler.
  \item \textbf{Feature-basierte Verfahren} suchen geometrische Punkte im
        Gesicht, zum Beispiel Augenwinkel oder die Mitte der Iris. Daraus
        berechnen sie den Blick.
  \item \textbf{Erscheinungsbasierte Verfahren} verwenden das Augen- oder
        Gesichtsbild selbst als Eingabe. Ein tiefes neuronales Netz sagt direkt
        2D-Koordinaten oder 3D-Winkel vorher.
\end{enumerate}
Durch den Fortschritt im Deep Learning sind erscheinungsbasierte Verfahren
heute der Standard unter \emph{Unconstrained Conditions} (dt.\
uneingeschränkte Bedingungen), also bei freier Kopfhaltung, wechselndem Licht
und beliebigem Abstand zur Kamera.

Cheng et al.~\cite{cheng2024review} teilen Verfahren mit \emph{Remote Camera}
(dt.\ entfernte Kamera: Die Kamera steht vor der Person und wird nicht am
Kopf getragen) ebenfalls in drei Gruppen ein: Verfahren mit einem
rekonstruierten 3D-Augenmodell (der Augapfel wird geometrisch nachgebildet,
und der Blick wird aus Pupille und Lichtreflexen berechnet), Regression auf
2D-Augen-Features und erscheinungsbasierte Verfahren. Deep-Learning-Ansätze ordnen sie nach vier Gesichtspunkten:
Feature Extraction, Netzarchitektur, persönliche Kalibrierung sowie Geräte und
Plattformen. Ein wichtiger Beitrag ist ein Benchmark. Er vereinheitlicht die
Vorverarbeitung (Gesichts- und Augenerkennung, Normalisierung der Daten) und
die Nachverarbeitung (Umrechnung zwischen 2D-Blickpunkt und 3D-Blickvektor).
So kann man Verfahren fair vergleichen. Bei kurzem Abstand zur Kamera oder
hoher Auflösung erreichen erscheinungsbasierte Verfahren laut Benchmark etwa
$5^\circ$. Bei großem Abstand und niedriger Auflösung (zum Beispiel Gaze360 und
RT-GENE) sind die Fehler deutlich größer. Als offene Themen nennen die Autoren:
(1) robustere Features, die nicht von der Umgebung abhängen, (2) hohe
Genauigkeit mit schneller und einfacher Kalibrierung -- hier gibt es einen
Zielkonflikt zwischen Kalibrierdauer und Genauigkeit -- und (3) die
Erklärbarkeit der gelernten Features.

Weitere Übersichtsarbeiten geben zusätzlichen Hintergrund. Liu et
al.~\cite{liu2022beholder} beschreiben den Weg von klassischen Verfahren zu
Deep Learning; ich nutze diese Arbeit vor allem als historische Referenz.
Pathirana et al.~\cite{pathirana2022survey} behandeln Deep-Learning-Verfahren
für einen und für mehrere Nutzer. Sie vergleichen CNN-Architekturen, von denen
einige auf ImageNet vortrainiert sind. Neuere vortrainierte Modelle wie DINOv2,
CLIP oder MAE stehen in diesen Übersichten aber nicht im Mittelpunkt. Für
Vergleiche mit Benchmarks nutze ich daher vor allem Cheng et
al.~\cite{cheng2024review}.

% ------------------------------------------------------------------------------
\section{Webcam-basiertes Eye-Tracking: Alltag und Labor}
\label{sec:webcam}
% ------------------------------------------------------------------------------
Webcams sind billig und fast überall vorhanden: in Laptops, an Desktop-Rechnern
und in Smartphones. Man braucht keine zusätzliche Hardware. Dafür sind die
Bedingungen viel schwieriger als bei speziellen Eye-Trackern mit Infrarotlicht:

\begin{itemize}
  \item \textbf{Infrarot-Systeme} nutzen hochauflösende Kameras, Infrarotlicht
        und Reflexe auf der Hornhaut. Liu et al.~\cite{liu2022beholder} nennen
        für solche Verfahren typische Genauigkeiten von $0,5^\circ$ bis
        $2^\circ$. Kommerzielle Geräte arbeiten in festen Abständen (zum
        Beispiel \SIrange{50}{80}{\centi\metre}) und erreichen unter optimalen
        Bedingungen $0,3^\circ$ bis $0,5^\circ$.
  \item \textbf{Webcams im Alltag} liefern Bilder mit niedriger Auflösung und
        ohne kontrolliertes Licht. Das Umgebungslicht ändert sich, die Lage
        von Kamera und Bildschirm ist bei jedem Gerät anders, und Abstand und
        Kopfhaltung sind frei. Liu et al.~\cite{liu2022beholder} nennen für
        erscheinungsbasierte Verfahren typische Fehler über $3^\circ$. Der
        Benchmark von Cheng et al.~\cite{cheng2024review} zeigt etwa $5^\circ$
        bei kurzem Abstand und deutlich mehr bei großem Abstand und niedriger
        Auflösung.
\end{itemize}

Zwei bekannte Arbeiten zeigen trotzdem, dass Webcam-Eye-Tracking praktisch
nutzbar ist. Papoutsaki et al.~\cite{papoutsaki2016webgazer} stellen
\emph{WebGazer} vor, einen Open-Source-Eye-Tracker, der komplett im Browser
läuft. WebGazer braucht keine eigene Kalibrierung. Er lernt mit
\emph{Ridge Regression} laufend aus Mausklicks und Mausbewegungen, welche
Augen-Features zu welcher Bildschirmposition gehören. Ridge Regression ist eine
lineare Regression mit einer Strafe für große Gewichte
(L2-Regularisierung); dadurch bleibt das Modell auch bei wenigen und
verrauschten Daten stabil. Die Genauigkeit reicht, um
grob den betrachteten Bereich zu erkennen, aber nicht für genaue Analysen.
Valliappan et al.~\cite{valliappan2020smartphone} trainieren dagegen ein CNN
auf GazeCapture vor, passen es mit Kalibrierdaten an und ergänzen ein kleines
Regressionsmodell pro Person. Auf Smartphones bei
\SIrange{25}{40}{\centi\metre} Abstand erreichen sie $0,6^\circ$ bis $1^\circ$.
Etwa 30~Sekunden Kalibrierung senken den Fehler von \SI{1,92}{\centi\metre} auf
\SI{0,46}{\centi\metre}. Laut den Autoren ist das so genau wie viel teurere
mobile Eye-Tracker. Die Bedingungen waren aber günstig (Kopf fast frontal,
kurzer Abstand). Beide Arbeiten zeigen: Ohne spezielle Hardware ist die
Personalisierung der wichtigste Faktor für die Genauigkeit.

Auch Eigenschaften der Person beeinflussen die Genauigkeit, zum Beispiel
Augenform, Gesichtsform, Brille und die gewohnte Kopfhaltung. Allgemeine
Modelle haben deshalb bei neuen Personen oft einen systematischen Versatz. Aus
diesem Grund ist die Personalisierung ein zentrales Forschungsthema
(\cref{sec:personalisierung}). Diese Arbeit untersucht genau solche
Bedingungen, unter anderem zwei Geräte mit verschiedenen Webcams, Brillen und
Hintergründe.

% ------------------------------------------------------------------------------
\section{Große Datensätze zur Blickschätzung}
\label{sec:datensaetze}
% ------------------------------------------------------------------------------
Moderne erscheinungsbasierte Verfahren brauchen große und vielfältige
Datensätze. \Cref{tab:datensaetze} fasst die für diese Arbeit wichtigen
Datensätze zusammen.

\begin{itemize}
  \item \textbf{MPIIGaze}~\cite{zhang2015mpiigaze} ist einer der ersten
        standardisierten Datensätze aus dem Alltag. Er enthält \num{213659}
        Gesichtsbilder von 15 Personen. Sie wurden über mehrere Monate bei der
        normalen Arbeit am Laptop aufgenommen~\cite{zhang2019mpiigaze}. Die
        Variante MPIIFaceGaze mit Gesichtsbildern wird meist mit
        Leave-one-subject-out ausgewertet (jede Person ist einmal die
        Testperson).
  \item \textbf{GazeCapture}~\cite{krafka2016itracker} enthält \num{2445504}
        Bilder von \num{1474} Personen, aufgenommen per Crowdsourcing auf
        iPhones und iPads. Licht und Kopfhaltung sind sehr verschieden.
        GazeCapture ist die Datengrundlage der Basismodelle dieser Arbeit.
        Gründe für diese Wahl: Es ist einer der größten und vielfältigsten
        Datensätze (viele Personen, Geräte und Umgebungen), und seine Labels
        sind 2D-Punkte auf dem Bildschirm (Point of Regard) -- also genau die
        Größe, die diese Arbeit schätzt.
  \item \textbf{ETH-XGaze}~\cite{zhang2020ethxgaze} ist ein hochauflösender
        Datensatz mit mehreren Kameras. Er enthält mehr als eine Million Bilder
        von 110 Personen mit sehr großen Kopf- und Blickwinkeln.
  \item \textbf{EVE}~\cite{park2020eve} enthält gleichzeitig aufgenommene
        Bildschirminhalte, Kamerabilder und Blickdaten. 54 Personen wurden mit
        vier Kameras aufgenommen, insgesamt mehr als 12~Millionen Bilder.
  \item \textbf{Gaze360}~\cite{kellnhofer2019gaze360} enthält Aufnahmen von
        238 Personen drinnen und draußen. Die 3D-Blicklabels decken sehr viele
        Kopfhaltungen und Abstände ab (etwa \num{172000} Bilder in Training,
        Validierung und Test).
  \item \textbf{EyeDiap} und \textbf{RT-GENE} werden vor allem für Tests über
        Datensätze hinweg verwendet. RT-GENE enthält unter anderem größere
        Kameraabstände und am Kopf getragene Geräte.
\end{itemize}

\begin{table}[htbp]
  \centering
  \small
  \caption{Überblick über ausgewählte Datensätze zur Blickschätzung.}
  \label{tab:datensaetze}
  \begin{tabularx}{\textwidth}{L{2.4cm} L{2.6cm} L{2.2cm} Y}
    \toprule
    \textbf{Datensatz} & \textbf{Umfang} & \textbf{Personen} & \textbf{Aufnahmebedingungen} \\
    \midrule
    MPIIGaze      & \num{213659} Bilder & 15 & Laptop-Webcam, Alltag über Monate \\
    GazeCapture   & \num{2445504} Bilder & \num{1474} & Smartphones und Tablets, Crowdsourcing \\
    ETH-XGaze     & >1~Million Bilder & 110 & Labor mit mehreren Kameras, extreme Kopfhaltungen \\
    EVE           & >12~Millionen Bilder & 54 & Bildschirm und Kameras, Videosequenzen \\
    Gaze360       & ca.\ \num{172000} Bilder & 238 & Drinnen und draußen, viele Kopfhaltungen und Abstände \\
    \bottomrule
  \end{tabularx}
\end{table}

% ------------------------------------------------------------------------------
\section{Netzarchitekturen für die Gaze Regression}
\label{sec:architekturen}
% ------------------------------------------------------------------------------
\subsection{Convolutional Neural Networks}
\label{subsec:cnn}
CNNs wie ResNet-18/50~\cite{he2016resnet}, VGG oder MobileNet sind die
üblichen Feature Extractors für Augen- und Gesichtsbilder. Ihre lokalen
Faltungen erkennen feine Texturen und Kanten in der Augenregion. Auch
spezielle Blickmodelle bauen oft auf CNNs auf: FAZE~\cite{park2019faze} nutzt
zum Beispiel einen Encoder-Decoder auf DenseNet-Basis, um Blick, Kopfhaltung
und Aussehen voneinander zu trennen. In dieser Arbeit ist ResNet-18 das
verwendete Modell.

\subsection{Vision Transformer}
\label{subsec:vit}
Vision Transformer (ViTs)~\cite{dosovitskiy2021vit} teilen ein Bild in kleine
Felder (Patches). Mit Self-Attention erfassen sie Zusammenhänge über das ganze
Bild, zum Beispiel zwischen Gesicht, Kopfdrehung und Augen.
UniGaze~\cite{qin2026unigaze} zeigt, dass mit MAE vortrainierte ViTs (Größen
Base, Large und Huge) sehr gute Backbones für die Blickschätzung sind.
TTP-GE~\cite{nagpure2025ttpge} kombiniert dagegen zwei ResNet-18-Zweige mit
einer Transformer-Fusion. GazeCLIP~\cite{wang2026gazeclip} zeigt, dass
Transformer nicht immer besser sind: Dort erreicht ein CLIP-vortrainiertes
ResNet-50 $3,5^\circ$ auf MPIIFaceGaze, die Variante mit ViT-B/32 nur
$3,7^\circ$. Ye et al.~\cite{ye2024transgaze} zeigen dagegen mit
\emph{TransGaze}, dass auch einfache ViTs mit großem Pretraining sehr gute
Ergebnisse auf GazeCapture und MPIIFaceGaze erreichen. Am besten war dort das
DINO-Pretraining mit \emph{Self-Supervised Learning}. Man sollte also CNNs und Transformer
gleichermaßen prüfen und nicht eine Architektur für allgemein besser halten.

\subsection{Vortrainierte Foundation-Modelle}
\label{subsec:foundation}
Beim Transfer Learning nutzt man ein Modell, das auf sehr vielen Bildern
vortrainiert wurde. Seine Features sind allgemein nutzbar. Man kann den
Backbone einfrieren und nur einen kleinen Regression Head darauf trainieren.
Für diese Arbeit sind vier Arten von Pretraining wichtig:
\begin{itemize}
  \item \textbf{Supervised Pretraining auf ImageNet} (dt.\ überwachtes
        Pretraining): Das Modell lernt mit Labels, die Menschen vergeben
        haben (hier: Objektklassen). Das ist der klassische Ausgangspunkt für
        CNNs wie ResNet~\cite{he2016resnet}. Diese Arbeit verwendet ein
        solches Modell.
  \item \textbf{DINOv2}~\cite{oquab2023dinov2}: lernt robuste Features durch
        \emph{Self-Supervised Learning}: Das
        Lernziel entsteht aus den Bildern selbst, ohne manuelle Labels.
  \item \textbf{CLIP}~\cite{radford2021clip}: lernt Bild und Text gemeinsam
        durch kontrastives Lernen auf Bild-Text-Paaren.
  \item \textbf{Masked Autoencoders (MAE)}~\cite{he2022mae}: lernen, indem sie
        verdeckte Bildteile rekonstruieren.
\end{itemize}
Wie man diese Features für die Blickschätzung nutzt, beschreibt der nächste
Abschnitt.

% ------------------------------------------------------------------------------
\section{Vortrainierte Modelle für die Blickschätzung}
\label{sec:pretrained}
% ------------------------------------------------------------------------------
\subsection{Eingefrorener Backbone oder vollständiges Fine-Tuning}
\label{subsec:frozen_full}
Es gibt zwei Grundstrategien, ein vortrainiertes Modell an die
Blickschätzung anzupassen. Beim \emph{vollständigen Fine-Tuning} werden alle
Gewichte trainiert. Das Modell passt sich so am stärksten an, braucht aber
viele Daten. Mit wenigen Daten droht Overfitting oder
\emph{katastrophales Vergessen}: Das Modell verliert das, was es im Pretraining
gelernt hat. Bei einem \emph{eingefrorenen Backbone} (zum Beispiel ResNet-18
oder ViT-B/14) bleiben die vortrainierten Gewichte fest. Man trainiert nur
einen kleinen Regression Head. So gibt es viel weniger trainierbare Parameter,
und die allgemeinen Features bleiben erhalten. Das ist wichtig, wenn man nur
wenige Kalibrierdaten hat.

Kumar et al.~\cite{kumar2022finetuning} begründen diesen Zielkonflikt
theoretisch und mit Experimenten. Sie vergleichen vollständiges Fine-Tuning mit
\emph{Linear Probing}, bei dem nur die letzte lineare Schicht trainiert wird.
Auf zehn Datensätzen mit Verteilungsverschiebung ist Fine-Tuning auf Daten wie
im Training im Mittel um 2\,\% besser, auf abweichenden Daten aber um 7\,\%
schlechter. Der Grund: Wenn Head und Backbone gleichzeitig trainiert werden und
der Head am Anfang noch zufällig ist, werden die vortrainierten Features
verzerrt. Die Autoren schlagen deshalb LP-FT vor: zuerst Linear Probing, dann
Fine-Tuning. Für die Personalisierung ist das direkt wichtig. Die
Kalibrierdaten einer Person sind wenige und decken die spätere Nutzung nur zum
Teil ab. Ein eingefrorener Backbone kann hier die Generalisierung schützen.

Der Regression Head ist meist ein mehrschichtiges Perzeptron (MLP). Es bildet
die Features, zum Beispiel die 512 Zahlen eines ResNet-18, auf die Koordinaten
$(x, y)$ ab. Wichtige Entscheidungen sind die Activation am Ausgang (zum
Beispiel Sigmoid, Gauß oder ReLU) und die Loss Function (zum Beispiel L1,
Smooth L1 oder Varianten mit Gewichtung nach der Position). Zwischen den beiden
Extremen liegen parametereffiziente Verfahren wie Adapter oder
LoRA~\cite{hu2022lora} (\cref{subsec:peft}).

\subsection{Großes MAE-Pretraining: UniGaze}
\label{subsec:unigaze}
Qin et al.~\cite{qin2026unigaze} stellen mit UniGaze ein großes Pretraining
speziell für die Blickschätzung vor. Mit Masked Autoencoders lernt das Modell
aus etwa 1,6~Millionen normalisierten Gesichtsbildern von mehr als
\num{260000} Personen. Die Daten kommen aus echten Quellen (VGGFace2, VFHQ,
CelebV-Text) und aus synthetischen Datensätzen (FaceSynthetics, SFHQ-T2I,
FFHQ-NV, XGaze-Dense). Ein zentrales Ergebnis ist, wie wichtig die
Normalisierung der Daten ist. Dabei werden Gesichtsbild und Blicklabel durch
Drehen und Skalieren in einen gemeinsamen Raum gebracht. So verschwindet ein
großer Teil der Unterschiede durch Kopfhaltung und Abstand. Zhang et
al.~\cite{zhang2018normalization} haben dieses Verfahren als Erste genau
untersucht. Eine Variante ohne Skalierungsfaktor verbessert die Genauigkeit je
nach Test um 9,5 bis 32,7\,\%. UniGaze bestätigt das: In den Ablationen
(UniGaze-B, nur echte Pretraining-Daten, Leave-one-dataset-out) senkt die
Normalisierung den Winkelfehler auf XTest von $14,56^\circ$ auf $11,95^\circ$.
Mit zusätzlichen synthetischen Daten werden $11,78^\circ$ erreicht.

UniGaze nutzt ViTs der Größen Base, Large und Huge. Beim Pretraining werden
75\,\% des Bildes verdeckt. ViT-Huge wird 300 Epochs lang vortrainiert. Das
dauert etwa \SI{120}{\hour} auf vier NVIDIA-H100-GPUs. Danach werden die
Modelle mit L1-Loss auf MPIIFaceGaze, ETH-XGaze, EYEDIAP, Gaze360 und
GazeCapture mit vollständigem Fine-Tuning angepasst. Die Autoren testen drei Protokolle:
\begin{itemize}
  \item \textbf{Domain Generalization} (Training auf ETH-XGaze oder Gaze360,
        Test auf MPIIFaceGaze oder dem Bildschirm-Teil von EYEDIAP):
        UniGaze-H erreicht $5,57^\circ$ (X$\rightarrow$M), $4,65^\circ$
        (X$\rightarrow$ECS), $5,65^\circ$ (G360$\rightarrow$M) und $5,86^\circ$
        (G360$\rightarrow$ECS). Das ist besser als andere Verfahren, zum
        Beispiel CLIP-Gaze mit $6,41^\circ$, $7,51^\circ$, $6,89^\circ$ und
        $7,06^\circ$.
  \item \textbf{Leave-one-dataset-out} (Training auf vier Datensätzen, Test
        auf dem fünften): UniGaze-H erreicht $11,29^\circ$ (XTest),
        $4,51^\circ$ (MPIIFaceGaze), $5,47^\circ$ (GazeCapture), $5,88^\circ$
        (EYEDIAP) und $12,37^\circ$ (Gaze360).
  \item \textbf{Gemeinsames Training auf allen Datensätzen}: UniGaze-H
        erreicht $4,46^\circ$ (XTest), $5,08^\circ$ (MPIIFaceGaze),
        $3,20^\circ$ (GazeCapture), $5,16^\circ$ (EYEDIAP) und $9,07^\circ$
        (Gaze360).
\end{itemize}

Die Ablationen zeigen: Mehr Pretraining-Daten verbessern die
Generalisierung, und echte plus synthetische Daten sind besser als nur echte
Daten. Auch Normalisierung und viele verschiedene Kopfhaltungen sind wichtig.
Wenn man die Vielfalt der Kopfhaltungen einschränkt, steigt der Fehler auf
XTest von $11,95^\circ$ auf $16,70^\circ$. Vergleiche mit DINO und FaRL deuten
darauf hin, dass die Zusammensetzung der Daten wichtiger sein kann als ihre
Menge.

\paragraph{Grenzen.}
UniGaze lässt einige Fragen offen, die für die Personalisierung mit Webcams
wichtig sind. Die Ergebnisse beruhen auf vollständigen Datensätzen und
vollständigem Fine-Tuning, nicht auf einer kurzen Kalibrierung mit wenigen
Bildern einer Person. Sparsame Strategien wie reines Head-Training, LoRA oder
Adapter stehen nicht im Mittelpunkt. Auch der Einfluss falscher Labels in
kurzen Kalibrierungen wird nicht untersucht. Das Pretraining braucht außerdem
sehr viel Rechenleistung (ViT-Huge: etwa \SI{120}{\hour} auf $4\times$H100).
Offen bleibt auch, ob die Verbesserungen unter echten Webcam-Bedingungen gelten
-- mit Brille, wechselndem Hintergrund, niedriger Auflösung und verschiedenen
Webcams.

\paragraph{Bedeutung für diese Arbeit.}
UniGaze ist eine wichtige Referenz für Pretraining speziell für den Blick.
UniGaze untersucht die Generalisierung über viele Personen. Diese Arbeit stellt
die ergänzende Frage: Wie gut kann man ein vortrainiertes Modell mit einer
kurzen Kalibrierung (hier \SIrange{30}{120}{\second}, also etwa 180 bis 720
Labels) an eine einzelne Person anpassen? Diese Arbeit verwendet dafür ein auf
ImageNet vortrainiertes ResNet-18. UniGaze-, DINOv2- und CLIP-Modelle sind ein
naheliegender nächster Schritt (\cref{tab:unigaze_vergleich}).

\begin{table}[htbp]
  \centering
  \small
  \caption{Vergleich von UniGaze und dieser Arbeit.}
  \label{tab:unigaze_vergleich}
  \begin{tabularx}{\textwidth}{L{2.8cm} L{3.2cm} L{2.6cm} Y}
    \toprule
    \textbf{Kriterium} & \textbf{UniGaze-H} & \textbf{UniGaze-B} & \textbf{Diese Arbeit} \\
    \midrule
    Pretraining-Daten & 1,6~Millionen Gesichtsbilder, >\num{260000} Personen & Gleiche Daten & ImageNet (ResNet-18), Regression Head auf GazeCapture \\
    Backbone & ViT-Huge & ViT-Base & ResNet-18, eingefroren \\
    Fine-Tuning & Gesamtes Modell & Gesamtes Modell & Nur Regression Head \\
    Kalibrierdaten & Vollständige Trainingsdatensätze & Vollständige Trainingsdatensätze & \SIrange{30}{120}{\second} pro Sitzung (ca.\ 180--720 Labels) \\
    Evaluation & Über Datensätze hinweg, Domain Gene\-rali\-zation & Über Datensätze hinweg, Domain Gene\-rali\-zation & Personendisjunkte Basismodelle, personalisierte Webcam-Sitzungen \\
    Falsche Labels & Keine Filterung nach Phase & Keine Filterung nach Phase & \code{WAIT}/\code{MOVE}-Protokoll \\
    Rechenaufwand & Ca.\ \SI{120}{\hour} auf $4\times$H100 & Nur für ViT-H angegeben & Kein Pretraining; Personalisierung ca.\ \SI{2}{\second} auf einer CPU \\
    Winkelfehler (Leave-one-dataset-out) & $4,51^\circ$--$12,37^\circ$ & $4,73^\circ$--$12,41^\circ$ & Diagonalfehler in \%, nicht direkt vergleichbar \\
    \bottomrule
  \end{tabularx}
\end{table}

\subsection{Bild-Text-Modelle: CLIP-Gaze und GazeCLIP}
\label{subsec:clip}
Zwei neuere Arbeiten nutzen das Bild-Text-Modell CLIP~\cite{radford2021clip}
für die Blickschätzung, aber mit verschiedenen Zielen.

\paragraph{CLIP-Gaze.}
Yin et al.~\cite{yin2024clipgaze} wollen mit CLIP-Gaze die
Domain Generalization verbessern. Das Modell soll in einer neuen Umgebung gut
bleiben, ohne gelabelte Daten aus dieser Umgebung. Die Idee: Faktoren, die
nichts mit dem Blick zu tun haben, werden durch Texte beschrieben, etwa
\enquote{An image of a face with \{condition\}.} Diese Faktoren bilden drei
Gruppen: Aussehen (Identität, Gesichtsausdruck, Haare, Alter usw.), Zubehör
(Brille, Hut, Make-up usw.) und Bildqualität (Noise, Unschärfe, Licht).
Drei Losses steuern das Training: Ein Distillation Loss richtet die
Blick-Features am CLIP-Image-Encoder aus. Ein Separation Loss macht sie den
blickfremden Features unähnlich. Ein Ranking Loss sorgt dafür, dass ähnliche
Blicke auch ähnliche Features haben. Zusätzlich passt die \emph{Personalized
Context Optimization} (PCO) die Texte an die jeweilige Person an. CLIP bleibt
eingefroren. Trainiert wird ein eigenes Netz aus ResNet-18, einem MLP-Filter
und einem Regression Head. Mit ETH-XGaze und Gaze360 als Quelle und
MPIIFaceGaze und EyeDiap als Ziel erreicht CLIP-Gaze über vier Aufgaben im
Mittel etwa $6,97^\circ$.

\paragraph{GazeCLIP.}
Wang et al.~\cite{wang2026gazeclip} nutzen grobe Blickinformation aus Text
für die Gaze Regression mit Supervised Learning. Ein Prompt-Modul bestimmt ohne Training
(Zero-Shot) mit Texten wie \enquote{A photo of a face gazing [class]} eine
grobe Richtung aus \{Front, Down, Left, Right\}. Der Image Encoder ist ein
CLIP-vortrainiertes ResNet-50 und wird mit Fine-Tuning weitertrainiert. Der Text Encoder (ein
Transformer mit 12 Schichten) bleibt eingefroren. Ein Cross-Attention-Modul
verbindet beide: Text ist die Query, die Bild-Features sind Key und Value. Ein
MLP mit drei Schichten ($1024 \rightarrow 256 \rightarrow 128 \rightarrow 2$)
sagt Pitch und Yaw vorher. GazeCLIP erreicht $3,5^\circ$ auf MPIIFaceGaze,
$4,7^\circ$ auf EyeDiap und $7,3^\circ$ auf RT-GENE. Cross-Attention
($3,6^\circ$) ist besser als einfaches Aneinanderhängen ($4,2^\circ$) oder
Addieren ($4,4^\circ$). Ein Image Encoder mit Fine-Tuning und eingefrorenem
Text Encoder ($3,6^\circ$) ist viel besser als zwei eingefrorene Encoder
($8,9^\circ$). Fehler entstehen vor allem bei schlechtem Licht und wenn Haare
oder Brillen die Augen verdecken. Außerdem kann die Zero-Shot-Klassifikation in
einer neuen Umgebung die falsche Richtung liefern.

\paragraph{Grenzen und Bedeutung für diese Arbeit.}
Beide Arbeiten zeigen, dass CLIP-Features für die Blickschätzung nützlich
sind. Eine Personalisierung mit kurzer Kalibrierung untersuchen sie aber nicht.
Bei CLIP-Gaze bleibt offen, ob eine Anpassung des CLIP-Image-Encoders die
Schätzung für eine Person verbessert. GazeCLIP braucht vollständiges
Supervised Training auf Benchmarkdaten. Es behandelt weder Personalisierung
mit wenigen Bildern noch falsche Kalibrierlabels, und die Daten gibt es laut
Veröffentlichung nur auf Anfrage. CLIP-Features wurden in dieser Arbeit noch
nicht untersucht; sie sind ein Kandidat für den Ausblick
(\cref{tab:clip_vergleich}). Die Winkelfehler dieser Arbeiten sind nur
Orientierungswerte. Sie lassen sich nicht direkt mit dem Diagonalfehler dieser
Arbeit vergleichen.

\begin{table}[htbp]
  \centering
  \small
  \caption{Vergleich von CLIP-Gaze, GazeCLIP und dieser Arbeit.}
  \label{tab:clip_vergleich}
  \begin{tabularx}{\textwidth}{L{2.6cm} Y Y Y}
    \toprule
    \textbf{Aspekt} & \textbf{CLIP-Gaze~\cite{yin2024clipgaze}} & \textbf{GazeCLIP~\cite{wang2026gazeclip}} & \textbf{Diese Arbeit} \\
    \midrule
    Ziel & Domain Gene\-rali\-zation & Supervised Gaze Regression & Personalisierung mit Webcam und wenigen Daten \\
    Nutzung von CLIP & Bild- und Text Encoder eingefroren & Image Encoder mit Fine-Tuning, Text Encoder eingefroren & Nicht verwendet (Ausblick) \\
    Architektur & ResNet-18, MLP-Filter, Regression Head & CLIP-ResNet-50, Cross-Attention, MLP & ResNet-18 (ImageNet), Regression Head \\
    Labels der Zielperson & Keine & Vollständige Trainingsdaten & \SIrange{30}{120}{\second} Kalibrierung \\
    Berichteter Fehler & Ca.\ $6,97^\circ$ (Mittel über 4 Aufgaben) & $3,5^\circ$ / $4,7^\circ$ / $7,3^\circ$ & Diagonalfehler in \% \\
    Datenqualität & Modelliert blickfremde Faktoren & Grobe Richtung aus Text (unsicher) & \code{WAIT}-Bilder werden entfernt \\
    \bottomrule
  \end{tabularx}
\end{table}

% ------------------------------------------------------------------------------
\section{Dateneffiziente Personalisierung}
\label{sec:personalisierung}
% ------------------------------------------------------------------------------
Allgemeine Modelle haben bei neuen Personen oft einen systematischen Versatz.
Er entsteht durch Unterschiede im Auge, in der Gesichtsform und in der
Kopfhaltung. Park et al.~\cite{park2019faze} sprechen von anatomischen
Unterschieden zwischen Personen, die die Genauigkeit allgemeiner Netze
begrenzen. Die Personalisierung soll diese Lücke mit möglichst wenig
Kalibrieraufwand schließen. Die folgenden Abschnitte beschreiben klassische
Kalibrierung, Few-Shot- und Meta-Learning sowie parametereffiziente Verfahren.

\subsection{Kalibrierverfahren}
\label{subsec:kalibrierung}
Bei der aktiven Kalibrierung schaut die Person nacheinander auf bekannte
Punkte auf dem Bildschirm. Die Koordinaten dieser Punkte sind die Labels für
die Anpassung. Typisch sind wenige feste Punkte oder ein Punkt, der sich über
den Bildschirm bewegt. iTracker~\cite{krafka2016itracker} verwendet bis zu 13
feste Punkte und trainiert darauf eine SVR pro Nutzer
(\cref{subsec:dl_pioniere}). He et al.~\cite{he2019ondevice} zeigen, dass
gute Verfahren mit nur 2 bis 5 Punkten pro Person besser sind als frühere
Verfahren mit mehr als 13 Punkten. Auf Smartphones ist ihr bestes Modell mit
höchstens fünf Punkten um 24 bis 26\,\% genauer als der damalige Stand der
Technik und läuft in Echtzeit auf dem Gerät. Sie schlagen außerdem ein
Verfahren vor, das nur ungelabelte Gesichtsbilder braucht. Valliappan et
al.~\cite{valliappan2020smartphone} verwenden neben Punkten eine
\emph{Zickzack-Kalibrierung}: Ein Ziel bewegt sich 60~Sekunden lang langsam
von links oben nach rechts unten, und die Frontkamera nimmt mit
\SI{30}{\hertz} auf. Schon etwa 100 Bilder senken den Fehler fast auf ein
Viertel.

Die Verfahren unterscheiden sich nicht nur darin, \emph{wie viele} Daten sie
brauchen, sondern auch darin, \emph{welchen Teil} des Modells sie anpassen:
\begin{itemize}
  \item \textbf{Offset-Kalibrierung:} Chen und Shi~\cite{chen2020offset}
        teilen den Blickwinkel in einen Teil, der für alle Personen gleich
        ist, und einen persönlichen Versatz (Bias). Sie beobachten, dass der
        Versatz zwischen Personen größer ist als die Schwankung bei einer
        Person. Die Kalibrierung braucht dann nur einen einzigen Blickpunkt.
        Auf MPIIGaze und EYEDIAP ist das Verfahren um bis zu 10\,\% besser als
        vergleichbare Modelle. Schon etwa neun Kalibrierbilder senken den
        Fehler deutlich.
  \item \textbf{Wenige Personenparameter:} Lindén et
        al.~\cite{linden2019spaze} beschreiben persönliche Unterschiede mit dem
        \emph{Spatial Adaptive Gaze Estimator} (SPAZE) durch wenige latente
        Parameter. Die Kalibrierung ist dann nur ein kleines
        Optimierungsproblem und kaum anfällig für Overfitting. Mit neun
        Beispielen erreicht SPAZE $2,70^\circ$ auf MPIIGaze, 14\,\% besser als
        der damalige Stand der Technik. Drei Parameter pro Auge reichen, um
        die persönlichen Unterschiede gut zu beschreiben.
  \item \textbf{Differenzielle Schätzung:} Liu et
        al.~\cite{liu2021differential} trainieren ein CNN, das den
        \emph{Unterschied} im Blick zwischen zwei Augenbildern derselben Person
        vorhersagt. Mit gelabelten Referenzbildern aus der Kalibrierung ergibt
        sich daraus der Blick im neuen Bild. Weil beide Bilder von derselben
        Person stammen, heben sich Störungen wie Augenform, Lidschluss oder
        Licht teilweise auf. Auf EYEDIAP, MPIIGaze und UT-Multiview ist das
        Verfahren schon mit einem einzigen Kalibrierbild besser als der Stand
        der Technik.
  \item \textbf{Personalisierung ohne Kalibrierung:} Zhao et
        al.~\cite{zhao2022ebpn} schlagen das \emph{Episode-based
        Personalization Network} (EbPN) vor. Ein eigener Zweig des Netzes soll
        persönliche Abweichungen ohne Kalibrierdaten ausgleichen.
\end{itemize}
Zusammen zeigen diese Arbeiten: Ein großer Teil des persönlichen Fehlers lässt
sich mit wenigen angepassten Parametern erklären. Das unterstützt die Strategie
dieser Arbeit, nur einen kleinen Regression Head auf einem eingefrorenen
Backbone anzupassen.

Cheng et al.~\cite{cheng2024review} nennen einen grundsätzlichen
Zielkonflikt: Längere Kalibrierungen machen die Schätzung genauer, sind aber
unbequemer. Hohe Genauigkeit mit kurzer, einfacher Kalibrierung ist daher eine
offene Forschungsfrage. Ein weiterer Punkt wird oft übersehen: die Qualität der
Kalibrierlabels selbst. Wenn ein Ziel springt, folgt der Blick erst nach einer
kurzen Verzögerung. Bilder aus dieser Transition Period können ein falsches Label
haben (\cref{sec:label_noise}).

\subsection{Few-Shot-Learning und Meta-Learning: FAZE}
\label{subsec:faze}
Park et al.~\cite{park2019faze} stellen mit FAZE (\emph{Few-Shot Adaptive Gaze
Estimation}) ein einflussreiches Verfahren für die Personalisierung mit wenigen
Kalibrierbildern vor. Es verbindet getrennte (entflochtene) Features für Blick,
Kopf und Aussehen mit Model-Agnostic Meta-Learning (MAML). FAZE trainiert nicht
für jede Person ein neues Modell. Es lernt stattdessen einen Startpunkt, von
dem aus wenige Beispiele für eine schnelle Anpassung reichen.

FAZE hat zwei Teile. Ein \emph{Disentangling Transforming Encoder-Decoder}
(DT-ED) auf DenseNet-Basis stellt Aussehen, Blickrichtung und Kopfhaltung in
getrennten Codes dar. Weil diese Codes mit Drehungen mitgehen
(rotationsäquivariant), kann man Änderungen von Blick und Kopf als Drehungen im
Feature Space beschreiben. Trainiert wird auf Bildpaaren derselben Person mit
einem Reconstruction Loss, einem Gaze Loss und einem Konsistenzterm. Der
zweite Teil ist ein Blickschätzer, der mit MAML meta-trainiert wird: Über viele
Trainingspersonen lernt er einen Startpunkt, aus dem wenige Gradientenschritte
ein persönliches Modell machen. Das soll Overfitting vermeiden, die beim
normalen Fine-Tuning mit sehr wenigen Daten droht.

FAZE wird auf GazeCapture trainiert und auf GazeCapture und MPIIGaze getestet.
Mit $k = 9$ Kalibrierbildern liegt der mittlere Fehler auf MPIIGaze bei
$3,14^\circ$. Für alle $1 \le k \le 256$ ist FAZE besser als die verglichenen
Verfahren. Auf GazeCapture sinkt der Fehler gegenüber dem Stand der Technik um
19\,\%, von $3,91^\circ$ auf $3,18^\circ$.

\paragraph{Grenzen und Abgrenzung.}
FAZE zeigt, dass eine Kalibrierung nicht nur von der Datenmenge abhängt,
sondern auch von den vorher gelernten Features und der Art der Anpassung. Die
Features von FAZE kommen aber aus Supervised Training mit Blickdaten und einem
DenseNet. Allgemeine vortrainierte Modelle (DINOv2, MAE, CLIP) werden nicht
untersucht. Die Anpassung nutzt Meta-Learning; einfachere Wege wie reines
Head-Training, LoRA oder Adapter stehen nicht im Mittelpunkt. Auch
Webcam-typische Probleme wie schwankende Bildqualität, Brillen, Hintergründe
und verschiedene Kameras werden nicht betrachtet, ebenso wenig das Entfernen von
Kalibrierbildern nach der Aufnahmephase. Diese Arbeit greift einige dieser
Punkte auf. Da FAZE Winkelfehler angibt und diese Arbeit einen Fehler relativ
zur Bildschirmdiagonale, sind die Zahlen nicht direkt vergleichbar.

\begin{table}[htbp]
  \centering
  \small
  \caption{Vergleich von FAZE und dieser Arbeit.}
  \label{tab:faze_vergleich}
  \begin{tabularx}{\textwidth}{L{3.0cm} Y Y}
    \toprule
    \textbf{Aspekt} & \textbf{FAZE~\cite{park2019faze}} & \textbf{Diese Arbeit} \\
    \midrule
    Ziel & Personalisierung der Blickrichtung mit wenigen Bildern & Dateneffiziente Personalisierung der Bildschirmposition mit Webcam \\
    Features & Entflochtene Codes, DenseNet-Encoder-Decoder & ResNet-18, auf ImageNet vortrainiert und eingefroren \\
    Pretraining & Supervised auf GazeCapture & ImageNet; Regression Head auf GazeCapture \\
    Anpassung & MAML und Gradientenschritte pro Person & Nur Regression Head \\
    Kalibrierdaten & $k = 1$ bis $256$ Bilder & \SIrange{30}{120}{\second} (ca.\ 180--720 Labels) \\
    Labelqualität & Keine Filterung nach Phase & \code{WAIT}-Bilder über Phasenprotokoll entfernt \\
    Evaluation & GazeCapture und MPIIGaze & 18 Webcam-Sitzungen unter wechselnden Bedingungen \\
    Berichteter Fehler & $3,14^\circ$ (MPIIGaze, $k = 9$) & Diagonalfehler in \%, nicht direkt vergleichbar \\
    \bottomrule
  \end{tabularx}
\end{table}

Yu et al.~\cite{yu2019redirection} erweitern die Anpassung mit wenigen
Bildern: Aus wenigen Referenzbildern einer Person erzeugen sie mit
Blickumlenkung (Gaze Redirection) zusätzliche Trainingsbilder und passen damit
ein allgemeines Modell an.

\subsection{Parametereffizientes Fine-Tuning (PEFT)}
\label{subsec:peft}
Parametereffiziente Verfahren passen nur einen kleinen Teil der Parameter an.
Der vortrainierte Backbone bleibt fast unverändert. Das spart Speicher und
Rechenzeit und senkt die Gefahr von Overfitting, wenn es nur wenige
Kalibrierdaten gibt. Für diese Arbeit sind drei Varianten wichtig:
\begin{itemize}
  \item \textbf{Reines Head-Training (nur FC):} Der Backbone bleibt komplett
        eingefroren, nur der Regression Head wird trainiert. Man kann die
        Features des Backbones einmal berechnen und speichern. Deshalb ist
        dieses Verfahren besonders schnell. Diese Arbeit verwendet es.
  \item \textbf{Adapter}~\cite{houlsby2019adapter}: Kleine zusätzliche Module
        werden zwischen die Schichten eines eingefrorenen Netzes gesetzt. Nur
        sie werden trainiert.
  \item \textbf{Low-Rank Adaptation (LoRA)}~\cite{hu2022lora}: Die
        vortrainierten Gewichtsmatrizen $W$ bleiben fest. Trainiert wird nur
        eine Änderung $\Delta W = BA$ aus zwei schmalen Matrizen mit niedrigem
        Rang $r$. So gibt es viel weniger trainierbare Parameter.
\end{itemize}
In der Blickschätzung wurden solche Verfahren bisher kaum systematisch
untersucht. Eine Ausnahme ist TPGaze~\cite{liu2024tpgaze}. Es überträgt das
Prompt-Tuning aus der Sprachverarbeitung: Ein zusätzlicher Prompt mit weniger
als 1\,\% der Parameter eines ResNet-18 wird zur Testzeit ohne Labels
angepasst. Meta-Learning sorgt dafür, dass diese Anpassung den Blickfehler
wirklich senkt (\cref{sec:tta}). Bao und Lu~\cite{bao2024feature2gaze} zeigen
außerdem, dass gerade die letzte vollständig verbundene Schicht, die viele
Bild-Features auf den Blick abbildet, leicht zu Overfitting neigt. Auch beim
reinen Head-Training ist der Aufbau des Heads also wichtig. Die anderen hier
betrachteten Arbeiten nutzen vollständiges Fine-Tuning (UniGaze, GazeCLIP),
Meta-Learning (FAZE) oder ein eingefrorenes Modell, das ein separat trainiertes
Netz steuert (CLIP-Gaze). Diese Arbeit untersucht reines Head-Training für die
Personalisierung mit kurzen Webcam-Kalibrierungen. Adapter und LoRA bleiben
für spätere Arbeiten.

% ------------------------------------------------------------------------------
\section{Domain Shift, Generalisierung und Evaluationsprotokolle}
\label{sec:domain_shift}
% ------------------------------------------------------------------------------
\subsection{Evaluation mit unbekannten Personen}
\label{subsec:norm_subject}
Wie gut ein Modell auf neue Personen generalisiert, hängt stark vom
Evaluationsprotokoll ab. Bei einer zufälligen Aufteilung
(\code{norm\_random}) können Bilder derselben Person in Training und Test
landen. In dieser Arbeit werden dabei Gruppen aus Person und Zielpunkt
aufgeteilt: Alle Bilder, in denen eine Person auf denselben Punkt schaut,
liegen zusammen entweder im Training oder im Test; dieselbe Person kommt aber
in beiden Teilen vor. Das Modell kann dann Features dieser Person ausnutzen, und die
Generalisierung wird überschätzt. Bei einem personendisjunkten Split
(\code{norm\_subject}) werden dagegen die Personen aufgeteilt (hier 80\,\%
Training, 20\,\% Test): Alle Bilder einer Person liegen nur im Training oder
nur im Test. Dieses Datenleck gibt es dann nicht (ähnlich wie bei
Leave-one-subject-out). Man misst dann die Leistung auf wirklich unbekannten Personen. In den
Basisexperimenten dieser Arbeit auf GazeCapture ergibt \code{norm\_subject}
einen Testfehler von \SI{28,39}{\percent} und \code{norm\_random}
\SI{27,72}{\percent}. Der größte Teil dieses Unterschieds kommt aber von den
unterschiedlichen Testdaten: Auch die Konstant-Baseline ist auf
\code{norm\_random} niedriger (\SI{27,66}{\percent} statt
\SI{28,46}{\percent}, siehe Kapitel~4).

Davon zu unterscheiden sind Tests über Datensätze hinweg, zum Beispiel
Leave-one-dataset-out in UniGaze~\cite{qin2026unigaze}. Dort wird ein ganzer
Datensatz zurückgehalten. Beide Protokolle messen also verschiedene Arten der
Generalisierung: auf unbekannte Personen in derselben Umgebung bzw.\ auf
unbekannte Aufnahmebedingungen.

\subsection{Domain Adaptation und Domain Generalization}
\label{subsec:domain_adaptation}
Wenn sich Trainings- und Testdaten in Personen, Kameras oder Umgebung
unterscheiden, werden erscheinungsbasierte Verfahren deutlich schlechter. Es
gibt zwei Gegenstrategien. Die Domain Adaptation nutzt (meist ungelabelte) Daten
aus der Zielumgebung. Die Domain Generalization braucht gar keine Zieldaten.

\paragraph{CRGA.}
Wang et al.~\cite{wang2022crga} schlagen mit \emph{Contrastive Regression Gaze
Adaptation} (CRGA) eine Unsupervised Domain Adaptation vor. Ein Modul für
\emph{Contrastive Domain Generalization} (CDG) lernt stabile Features in der
Source Domain: Features mit ähnlichem Blick werden zusammengezogen, Features mit
sehr unterschiedlichem Blick auseinandergeschoben. Ein Modul für
\emph{Contrastive Self-training Adaptation} (CSA) lernt danach mit
Pseudo-Labels in der Target Domain. Mit ETH-XGaze und Gaze360 als Quelle und
MPIIGaze, RT-GENE, GazeCapture und EyeDiap als Ziel verbessern die Autoren die
Baseline um 40,4\,\%, 34,7\,\%, 55,8\,\% und 34,3\,\% (Quelle ETH-XGaze) bzw.\
um 31,7\,\%, 30,5\,\%, 32,9\,\% und 23,8\,\% (Quelle Gaze360).

CRGA löst also ein anderes Problem: die Anpassung an eine neue Umgebung ohne
gelabelte Zieldaten. Diese Arbeit nutzt dagegen wenige gelabelte
Kalibrierdaten der Zielperson. Damit ist Supervised Training mit echten Labels
möglich, nicht nur mit Pseudo-Labels. Selbsttraining mit Pseudo-Labels kann
außerdem einen Bestätigungsfehler (Confirmation Bias) erzeugen: Falsche
Vorhersagen werden im weiteren Training verstärkt. Zudem geht es hier um
Unterschiede zwischen Personen und nicht um den Wechsel zwischen
Benchmark-Datensätzen. CRGA ist daher verwandte Arbeit zur Robustheit, aber
keine direkt vergleichbare Baseline.

\paragraph{Domain Generalization.}
Verfahren wie PureGaze~\cite{cheng2022puregaze} oder
CLIP-Gaze~\cite{yin2024clipgaze} (\cref{subsec:clip}) versuchen, blickfremde
Faktoren schon im Training aus den Features zu entfernen. Dann funktioniert
das Modell auch ohne Zieldaten in neuen Umgebungen.
UniGaze~\cite{qin2026unigaze} zeigt, dass großes blickspezifisches Pretraining
diese Verfahren bei Tests über Datensätze hinweg übertreffen kann
(\cref{subsec:unigaze}).

\subsection{Blickschätzung für mehrere Personen}
\label{subsec:gazeonce}
Zhang et al.~\cite{zhang2022gazeonce} stellen mit GazeOnce ein einstufiges
Verfahren vor. Es schätzt den Blick aller Personen in einem Bild in einem
einzigen Durchlauf und bleibt auch bei mehr als zehn Gesichtern
echtzeitfähig. Neben der Blickrichtung sagt es auch Gesichtsrahmen
(Bounding Boxes) und Gesichtspunkte (Landmarken) vorher. Ein Self-Supervised
Loss prüft, ob die geschätzte 3D-Richtung zu ihren 2D-Projektionen passt.
Für das Training erzeugen die Autoren den synthetischen Datensatz MPSGaze: Sie
setzen Gesichter aus ETH-XGaze mit bekanntem Blick in passende Gesichter aus
WiderFace ein. Laut den Autoren ist GazeOnce genauer und schneller als
bestehende Verfahren, besonders bei vielen Gesichtern.

GazeOnce löst damit ein anderes Problem: den Blick vieler Personen in freien
Szenen, nicht die Personalisierung für einen Webcam-Nutzer. Methodisch
interessant ist die Konsistenz zwischen 3D-Richtung und 2D-Projektion. MPSGaze
zeigt außerdem, dass synthetische Trainingsdaten sich von echten
Webcam-Bildern unterscheiden können.

\subsection{Unterschiede bei Geräten und Umgebung}
\label{subsec:geraete}
Neben Unterschieden zwischen Datensätzen beeinflussen viele Faktoren von Gerät
und Umgebung die Genauigkeit: Brille und Verdeckungen, Hintergrund und Licht,
Bildauflösung und Kameraabstand sowie die Lage von Kamera und Bildschirm. Die
Literatur gibt einige Hinweise:
\begin{itemize}
  \item Cheng et al.~\cite{cheng2024review} beobachten deutlich größere Fehler
        bei großem Abstand und niedriger Auflösung (Gaze360, RT-GENE).
  \item GazeCLIP~\cite{wang2026gazeclip} nennt schlechtes Licht und
        Verdeckung durch Haare oder Brille als typische Fehlerfälle.
  \item UniGaze~\cite{qin2026unigaze} betont, wie wichtig viele verschiedene
        Kopfhaltungen und ein großer Bereich der Blicklabels sind.
  \item CLIP-Gaze~\cite{yin2024clipgaze} modelliert Zubehör und Bildqualität
        ausdrücklich als blickfremde Faktoren.
\end{itemize}
Systematische Untersuchungen dieser Faktoren bei der Personalisierung mit
kurzen Webcam-Kalibrierungen fehlen aber weitgehend. Diese Arbeit betrachtet
deshalb über mehrere Sitzungen auch zwei Geräte mit verschiedenen Webcams,
Brille, Hintergrund und Kalibrierdauer (siehe Kapitel~4).

% ------------------------------------------------------------------------------
\section{Test-Time Adaptation}
\label{sec:tta}
% ------------------------------------------------------------------------------
Bei der \emph{Test-Time Adaptation} (TTA, dt.\ Anpassung zur Testzeit) wird ein fertig
trainiertes Modell erst während der Nutzung an die Person oder die Umgebung
angepasst, meist mit wenigen oder gar keinen Labels. Anders als bei der
klassischen Kalibrierung (\cref{subsec:kalibrierung}) soll der Nutzer dabei
möglichst wenig selbst tun müssen.

\subsection{TTP-GE: Self-Supervised und Weakly Supervised Adaptation}
\label{subsec:ttpge}
Nagpure et al.~\cite{nagpure2025ttpge} stellen mit TTP-GE (\emph{Test-Time
Personalized Gaze Estimation}) ein zweistufiges Verfahren für die
Blickschätzung im Auto vor. Es baut auf GazeDPTR~V2 auf. Diese Architektur hat
zwei Zweige: einen für das Originalbild des Gesichts und einen für ein
normalisiertes Gesichtsbild. Ein Transformer führt die Features zusammen.
Mit Positionsinformation aus einer Punktprojektion kann das Modell zusätzlich
Blickzonen erkennen.

Das Meta-Training nutzt zwei Schleifen auf Basis von MAML. In der inneren
Schleife passt ein Self-Supervised Consistency Loss (im Stil von BYOL++) das
Modell an neue Personen an -- ohne Blicklabels. Die äußere Schleife optimiert
Gaze Regression und Zonenerkennung mit gelabelten Daten. Zur Testzeit gibt es
zwei Schritte: Zuerst passt sich das Modell mit Self-Supervised Learning an ungelabelte
Bilder der Person an; das hilft bei veränderten Bildbedingungen. Danach wird
die Zonenerkennung mit wenigen Zonenlabels angepasst (zum Beispiel Spiegel,
Armaturenbrett, Windschutzscheibe). Diese Labels werden in einer kurzen
Einrichtungsphase gesammelt.

Auf IVGaze (\num{44705} Bilder von 125 Personen) erreicht die Kombination
einen mittleren Fehler von $6,00^\circ$. Die Baseline liegt bei $6,78^\circ$,
nur \emph{Weakly Supervised Adaptation} (dt.\ schwach überwachte Anpassung,
hier mit wenigen groben Zonenlabels) bei $6,13^\circ$ und nur
\emph{Self-Supervised Adaptation} (dt.\ selbstüberwachte Anpassung, ganz ohne
Labels) bei $6,72^\circ$. Der Anteil der Schätzungen mit weniger als
$6^\circ$ Fehler steigt von 49,95\,\% auf 58,2\,\%. Die beiden Schritte
ergänzen sich also: Die Selbstüberwachung behandelt die Bildbedingungen, die
Zonenlabels liefern persönliche Information.

\subsection{Weitere labelsparende Verfahren}
\label{subsec:tta_weitere}
Einige neuere Arbeiten verfolgen ähnliche Ziele mit anderen Mitteln:
\begin{itemize}
  \item \textbf{TPGaze}~\cite{liu2024tpgaze} passt zur Testzeit nur einen
        meta-gelernten Prompt an. Er hat weniger als 1\,\% der Parameter eines
        ResNet-18. Als Ziel ohne Labels reicht ein einfacher Symmetry Loss.
        Laut den Autoren ist die Anpassung bis zu zehnmal schneller als bei
        anderen Verfahren. Getestet wird in vier Szenarien über Datensätze
        hinweg.
  \item \textbf{ELF-UA}~\cite{wu2024elfua} behandelt die Anpassung
        \emph{ohne Labels} als eigenes Problem: Wenige ungelabelte Bilder der
        Person reichen. Das Meta-Learning nutzt eine Schranke aus der
        Domain Adaptation als Loss Function.
  \item \textbf{PIGaze}~\cite{wang2026pigaze} gewinnt Kalibrierdaten im Auto
        automatisch: Der Blick auf bekannte Zonen wie Spiegel und
        Windschutzscheibe muss geometrisch zusammenpassen
        (Selbstkalibrierung). Ein vortrainiertes Basismodell ist ein
        \emph{eingefrorener} Feature Extractor. Darauf lernt ein kleiner
        Adapter die persönliche Abbildung. Auf IVGaze und LBW ist PIGaze
        besser als Fine-Tuning.
\end{itemize}

\paragraph{Grenzen und Abgrenzung.}
Diese Verfahren wollen den Kalibrieraufwand mit ungelabelten oder schwach
gelabelten Daten senken. Sie schätzen meist 3D-Blickrichtungen und arbeiten bei
TTP-GE und PIGaze im Auto mit festen Zonen. Ihre Anpassung (Meta-Training,
Online-Updates, EMA-Zielnetze) ist deutlich aufwendiger als einfaches
Head-Training. Allgemeine vortrainierte Modelle wie DINOv2, CLIP oder MAE werden
nicht als Backbone untersucht. Diese Arbeit verfolgt den ergänzenden Weg einer
\emph{Supervised} Personalisierung mit Bildschirmkoordinaten aus einer kurzen
Kalibrierung. Interessant ist aber, dass PIGaze -- wie diese Arbeit -- einen
eingefrorenen Backbone mit einem kleinen, persönlich trainierten Modul
kombiniert.

\begin{table}[htbp]
  \centering
  \small
  \caption{Vergleich von TTP-GE und dieser Arbeit.}
  \label{tab:ttpge_vergleich}
  \begin{tabularx}{\textwidth}{L{3.0cm} Y Y}
    \toprule
    \textbf{Aspekt} & \textbf{TTP-GE~\cite{nagpure2025ttpge}} & \textbf{Diese Arbeit} \\
    \midrule
    Ansatz & Test-Time Adaptation mit ungelabelten Bildern und wenigen Zonenlabels & Supervised Personalisierung mit gelabelten Kalibrierdaten \\
    Backbone & GazeDPTR~V2: zwei ResNet-18-Zweige mit Transformer-Fusion & ResNet-18 (ImageNet), eingefroren \\
    Labels & Ungelabelte Bilder und neun Blickzonen & Bildschirmkoordinaten; \code{WAIT}-Bilder werden entfernt \\
    Anpassung & Meta-Training, Online-Updates, EMA-Zielnetz & Nur Regression Head \\
    Anwendung & Fahrerüberwachung im Auto & Webcam vor einem Bildschirm \\
    Ergebnis & $6,00^\circ$ auf IVGaze & Diagonalfehler in \%, nicht direkt vergleichbar \\
    \bottomrule
  \end{tabularx}
\end{table}

% ------------------------------------------------------------------------------
\section{Datenqualität, falsche Labels und zeitlicher Ablauf}
\label{sec:label_noise}
% ------------------------------------------------------------------------------
Supervised Personalisierung setzt voraus, dass jedes Kalibrierbild das
richtige Label hat, also den Punkt, auf den die Person in diesem Moment
wirklich schaut. In der Praxis stimmt das nur ungefähr. Abweichungen zwischen
Label und echtem Blick nennt man \emph{Label Noise} (dt.\ Label-Rauschen). Bei kurzen
Kalibrierungen ist es besonders schädlich, weil jedes falsche Bild einen
großen Anteil der Trainingsdaten ausmacht.

\subsection{Datenqualität im Eye-Tracking}
\label{subsec:datenqualitaet}
In der Eye-Tracking-Forschung beschreibt man die Datenqualität meist mit drei
Größen. Niehorster et al.~\cite{niehorster2026quality} geben einen aktuellen
Überblick, wie man sie misst, und stellen Werkzeuge dafür bereit. Für $N$
Messungen, bei denen die Person auf ein bekanntes Ziel schaut, gilt:
\begin{itemize}
  \item \textbf{Accuracy} (dt.\ Genauigkeit): die \emph{systematische}
        Abweichung zwischen gemessenem und wahrem Blickpunkt, meist als
        mittlerer Winkelfehler
        \[
          \mathrm{Acc} = \frac{1}{N}\sum_{i=1}^{N} \theta_i ,
        \]
        wobei $\theta_i$ der Winkel zwischen gemessenem und wahrem Blick ist.
  \item \textbf{Precision} (dt.\ Präzision): die \emph{Streuung} wiederholter
        Messungen, unabhängig vom wahren Ziel. Üblich ist der
        RMS-Wert zwischen aufeinanderfolgenden Messungen
        \[
          \mathrm{RMS\text{-}S2S} = \sqrt{\frac{1}{N-1}\sum_{i=1}^{N-1} \theta_{i,i+1}^{2}} ,
        \]
        mit $\theta_{i,i+1}$ als Winkel zwischen Messung $i$ und $i+1$, oder
        die Standardabweichung der Messungen um ihren Mittelwert.
  \item \textbf{Data Loss} (dt.\ Datenverlust): der Anteil der Messungen
        ohne gültigen Blickwert,
        $\mathrm{DL} = N_{\text{ungültig}} / N \cdot \SI{100}{\percent}$.
\end{itemize}
Der Unterschied: Ein Eye-Tracker kann präzise, aber ungenau sein (die
Messungen liegen dicht beieinander, aber neben dem Ziel), oder genau, aber
unpräzise (der Mittelwert stimmt, aber die Messungen streuen stark). Der
Diagonalfehler dieser Arbeit ist ein RMSE und enthält damit beides: die
systematische Abweichung und die Streuung. Bei
lernenden Verfahren kommt eine vierte Größe dazu: die Qualität der
\emph{Trainingslabels}. Falsche Labels verfälschen nicht nur die Auswertung,
sondern auch das gelernte Modell.

\subsection{Reaktionszeit und Saccades}
\label{subsec:latenz}
Wenn ein Kalibrierziel an einer neuen Stelle erscheint, folgt der Blick nicht
sofort. Bis die \emph{Saccade} beginnt (dt.\ Saccade: eine sehr schnelle
Augenbewegung zu einem neuen Ziel), vergeht eine kurze Reaktionszeit. Valliappan et al.~\cite{valliappan2020smartphone} messen auf
Smartphones im Mittel \SI{210}{\milli\second} (Median \SI{167}{\milli\second}).
Das passt zu früheren Studien mit \SIrange{200}{250}{\milli\second}. Danach
kommen die Saccade selbst und eventuell kleine Korrekturen, bis das Auge ruhig
auf dem Ziel bleibt.

Bilder aus dieser \emph{Transition Period} (dt.\ Übergangszeit) haben das
Label der \emph{neuen} Position. Der
Blick ist aber noch auf der \emph{alten} Position oder irgendwo dazwischen. Bei
\SI{30}{fps} betrifft allein eine Reaktionszeit von etwa
\SI{200}{\milli\second} rund sechs Bilder pro Sprung. Je öfter das Ziel
springt, desto mehr Bilder haben ein falsches Label.

\subsection{Filtern schon bei der Aufnahme}
\label{subsec:phase_logging}
Eine einfache Gegenmaßnahme: Man verwirft die Bilder direkt nach einem
Zielwechsel. Valliappan et al.~\cite{valliappan2020smartphone} verwerfen zum
Beispiel immer die ersten \SI{800}{\milli\second} nach jedem neuen Ziel. Ein
festes Zeitfenster ist leicht umzusetzen. Es beruht aber auf einer Annahme über
die Dauer des Übergangs, und seine Wirkung wird meist nicht eigens untersucht.

Diese Arbeit geht einen ähnlichen, aber genaueren Weg. Das Aufnahmeprogramm
speichert für jedes Bild die aktuelle \emph{Phase} des Ziels: \code{WAIT} (eine
kurze Wartezeit von \SI{0,5}{\second} direkt nach einem Sprung) oder
\code{MOVE} (das Ziel bewegt sich danach \SI{4,5}{\second} lang). So kann man
die \code{WAIT}-Bilder eindeutig und reproduzierbar erkennen und vor dem
Training entfernen. In den Aufnahmen betrifft das etwa 10\,\% der Bilder.
Anders als ein festes Zeitfenster erlaubt dieses Protokoll einen direkten
Vergleich: Modelle mit und ohne bereinigte Daten unter sonst gleichen
Bedingungen.

\subsection{Filtern nach der Aufnahme und robuste Lernverfahren}
\label{subsec:dmrf}
Andere Verfahren erkennen falsche Labels erst \emph{nach} der Aufnahme mit
statistischen oder modellbasierten Kriterien.

\paragraph{DMRF.}
Jiang et al.~\cite{jiang2026dmrf} schlagen mit \emph{Dynamic Multi-model
Regression Filtering} (DMRF) ein allgemeines Verfahren gegen falsche Labels in
Regressionsaufgaben vor. Es löst ein \emph{Chicken-and-Egg Problem} (dt.\ Henne-Ei-Problem: Zwei
Dinge hängen gegenseitig voneinander ab, sodass unklar ist, womit man
beginnen soll): Falsche Labels verzerren
die Modelle, mit denen man das Noise schätzt, und verzerrte Modelle schätzen
das Noise falsch. DMRF hat drei Teile:
\begin{enumerate}
  \item eine \emph{Aufwärmfilterung} mit einer Verteilungsfunktion der
        geschätzten Fehler,
  \item eine \emph{dynamische Modellauswahl} nach einem
        Korrelationskriterium,
  \item einen \emph{Rückhol-Schritt} (Data Recall), der Beispiele mit kleinem
        Loss wieder zu den sauberen Daten nimmt.
\end{enumerate}
Neben 15 Benchmarks aus UCI und KEEL und einer Altersschätzung wird DMRF in das
Blickverfahren L2CS-Net eingebaut. Auf einem vorverarbeiteten
MPIIGaze-Datensatz mit \num{45000} Bildern (Leave-one-person-out) sinkt der
mittlere Fehler von $3,92^\circ$ auf $3,85^\circ$. Dabei bleiben 74,3\,\% der
Daten erhalten. Der Code ist öffentlich.

\paragraph{SeeTN.}
Peng et al.~\cite{peng2026seetn} untersuchen als Erste systematisch, wie
falsche Labels die Domain Generalization in der Blickschätzung beeinflussen.
In einem Vorversuch verfälschen sie 10, 20 und 30\,\% der Gaze360-Labels mit
starkem Gaussian Noise. Der Fehler einer ResNet-18-Baseline auf MPIIGaze steigt
dadurch von $7,94^\circ$ auf $12,17^\circ$ (20\,\%) bzw.\ $14,38^\circ$
(30\,\%). Ihr Verfahren \emph{See-Through-Noise} baut einen semantischen
Feature Space auf. Es erkennt falsche Beispiele daran, dass Feature und Label
nicht zusammenpassen, und überträgt Blickinformation von sauberen auf falsche
Beispiele.

\paragraph{Vorher oder nachher filtern?}
Verfahren wie DMRF und SeeTN erkennen falsche Labels aus den Daten selbst. Sie
finden daher auch unerwartete Fehlerquellen. Sie brauchen aber genug Beispiele
für eine sichere Schätzung, rechnen aufwendiger und wurden bisher nur auf
großen Datensätzen getestet, nicht auf kurzen Kalibrierungen einer Person. Das
Phasenprotokoll erfasst dagegen nur eine \emph{bekannte} Fehlerquelle, die
durch den Ablauf der Aufnahme entsteht -- dafür aber eindeutig und ohne
zusätzliches Modell. Beide Ansätze ergänzen sich (\cref{tab:dmrf_vergleich}).

\begin{table}[htbp]
  \centering
  \small
  \caption{Vergleich von nachträglicher Filterung (DMRF) und Phasenprotokoll.}
  \label{tab:dmrf_vergleich}
  \begin{tabularx}{\textwidth}{L{3.0cm} Y Y}
    \toprule
    \textbf{Kriterium} & \textbf{DMRF~\cite{jiang2026dmrf}} & \textbf{Diese Arbeit} \\
    \midrule
    Strategie & Nachher: Fehler schätzen und filtern & Vorher: Phase bei der Aufnahme speichern \\
    Erkennung & Modellvorhersagen, Fehlermuster, dynamische Modellauswahl & \code{WAIT}/\code{MOVE}-Label pro Bild \\
    Erfasste Fehler & Alle statistisch auffälligen Beispiele & Nur die Transition Period nach einem Sprung \\
    Erhaltene Daten & 74,3\,\% (MPIIGaze) & Ca.\ 90\,\% \\
    Anwendung & Viele Personen, Leave-one-person-out & Webcam-Kalibrierung einer Person \\
    Zusatzaufwand & Mehrere Regressionsmodelle, iteratives Training & Kein zusätzliches Modell \\
    Ergebnis & $3,92^\circ \rightarrow 3,85^\circ$ & Diagonalfehler in \%, nicht direkt vergleichbar \\
    \bottomrule
  \end{tabularx}
\end{table}

% ------------------------------------------------------------------------------
\section{Vergleich und kritische Diskussion}
\label{sec:vergleich}
% ------------------------------------------------------------------------------
\Cref{tab:sota_vergleich} fasst die wichtigsten Arbeiten dieses Kapitels
zusammen. Die Spalten sind: Ziel, verwendete Features, Art der
Personalisierung und Umgang mit der Labelqualität.

\begin{table}[htbp]
  \centering
  \footnotesize
  \setlength{\tabcolsep}{3pt}
  \caption[Vergleich der wichtigsten Arbeiten]{Vergleich der wichtigsten Arbeiten. $k$ = Anzahl der Kalibrierbilder pro Person. Die Fehler stammen aus verschiedenen Protokollen und Metriken und sind nicht direkt vergleichbar.}
  \label{tab:sota_vergleich}
  \begin{tabularx}{\textwidth}{L{2.3cm} L{2.4cm} L{2.6cm} L{2.6cm} Y L{1.9cm}}
    \toprule
    \textbf{Arbeit} & \textbf{Ziel} & \textbf{Features} & \textbf{Personalisierung} & \textbf{Labelqualität} & \textbf{Ergebnis} \\
    \midrule
    iTracker~\cite{krafka2016itracker} & 2D-Blick, mobil & CNN-Zweige ähnlich AlexNet & SVR, bis zu 13 Punkte & -- & \SI{1,34}{\centi\metre} (Smartphone) \\
    FAZE~\cite{park2019faze} & Wenige Bilder & DenseNet-DT-ED, supervised & MAML, $k \le 256$ & -- & $3,14^\circ$ ($k=9$) \\
    SPAZE~\cite{linden2019spaze} & Wenige Bilder & CNN, latente Personenparameter & Kleines Optimierungsproblem & -- & $2,70^\circ$ ($k=9$) \\
    Chen \& Shi~\cite{chen2020offset} & Offset-Kalibrierung & CNN & Versatz aus einem Blickpunkt & -- & bis zu 10\,\% besser \\
    Valliappan~\cite{valliappan2020smartphone} & Smartphone & CNN, vortrainiert auf GazeCapture & Fine-Tuning + Regression, ca.\ 30\,s & Feste 800-ms-Verwerfung & \SI{0,46}{\centi\metre} \\
    CRGA~\cite{wang2022crga} & Unsupervised DA & ResNet (ImageNet) & Pseudo-Labels & -- & bis zu 55,8\,\% besser \\
    UniGaze~\cite{qin2026unigaze} & Generalisierung & MAE-ViT, 1,6~Millionen Gesichter & Vollständiges Fine-Tuning & -- & $4,65$--$5,86^\circ$ (DG) \\
    CLIP-Gaze~\cite{yin2024clipgaze} & Generalisierung & CLIP (eingefroren) + ResNet-18 & -- & Blickfremde Faktoren & $6,97^\circ$ (DG) \\
    GazeCLIP~\cite{wang2026gazeclip} & Supervised Regression & CLIP-ResNet-50 (Fine-Tuning) & -- & -- & $3,5^\circ$ (MPIIFaceGaze) \\
    TTP-GE~\cite{nagpure2025ttpge} & Test-Time Adaptation & GazeDPTR~V2 & Self- + weakly supervised & -- & $6,00^\circ$ (IVGaze) \\
    TPGaze~\cite{liu2024tpgaze} & Test-Time Adaptation & ResNet-18 & Prompt (<1\,\% der Parameter) & -- & -- \\
    DMRF~\cite{jiang2026dmrf} & Falsche Labels filtern & L2CS-Net & -- & Nachher, modellbasiert & $3,85^\circ$ (MPIIGaze) \\
    SeeTN~\cite{peng2026seetn} & Robuste Generalisierung & ResNet-18/50 & -- & Nachher, feature-basiert & -- \\
    \midrule
    \textbf{Diese Arbeit} & Webcam-Personalisierung & ResNet-18 (ImageNet), eingefroren & Nur Head; \SIrange{30}{120}{\second} & Vorher, Phasenprotokoll & Diagonalfehler in \% \\
    \bottomrule
  \end{tabularx}
\end{table}

\paragraph{Kritische Diskussion.}
Aus dem Vergleich ergeben sich vier Beobachtungen:
\begin{enumerate}
  \item \textbf{Features und Personalisierung werden getrennt untersucht.}
        Arbeiten zu großen vortrainierten Modellen (UniGaze, CLIP-Gaze,
        GazeCLIP) zielen auf Generalisierung über Datensätze hinweg und nutzen
        vollständiges Fine-Tuning mit vielen Daten. Arbeiten zur
        Personalisierung (FAZE, SPAZE, Chen \& Shi) nutzen dagegen CNNs, die
        speziell für den Blick trainiert wurden. Wie gut sich allgemeine
        vortrainierte Modelle mit wenigen Kalibrierdaten personalisieren
        lassen, ist offen.
  \item \textbf{Wenige Parameter reichen oft.} Offset-Kalibrierung, wenige
        Personenparameter, Prompt-Tuning und eingefrorene Modelle mit Adaptern
        (PIGaze) zeigen übereinstimmend: Ein großer Teil des persönlichen
        Fehlers lässt sich mit wenigen angepassten Parametern erklären. Kumar
        et al.~\cite{kumar2022finetuning} liefern dafür eine allgemeine
        Begründung.
  \item \textbf{Falsche Labels werden meist erst nachher behandelt.} DMRF und
        SeeTN filtern oder gewichten falsche Labels in großen Datensätzen.
        Fehler durch den Ablauf kurzer Kalibrierungen werden, wenn überhaupt,
        mit festen Zeitfenstern behandelt (Valliappan et al.), ohne ihren
        Einfluss genau zu untersuchen.
  \item \textbf{Die Ergebnisse sind kaum vergleichbar.} Die Arbeiten geben
        Winkel in Grad, Abstände in Zentimetern oder relative Verbesserungen an
        und nutzen verschiedene Datensätze, Splits und Protokolle. Einen
        Winkelfehler kann man nur mit bekannter Geometrie (Abstand,
        Bildschirmgröße) in einen Bildschirmfehler umrechnen. Die Metrik dieser
        Arbeit, der euklidische Fehler relativ zur Bildschirmdiagonale, ist
        daher nicht direkt mit der Literatur vergleichbar. Die Zahlen in
        \cref{tab:sota_vergleich} dienen der Orientierung, nicht als Rangliste.
\end{enumerate}

% ------------------------------------------------------------------------------
\section{Forschungslücken und Bezug zu dieser Arbeit}
\label{sec:luecken}
% ------------------------------------------------------------------------------
\subsection{Forschungslücken}
\label{subsec:luecken}
Aus dem Stand der Technik ergeben sich vier Forschungslücken:

\begin{enumerate}[label=\textbf{F\arabic*:}, leftmargin=*, align=left]
  \item \textbf{Dateneffiziente Personalisierung vortrainierter Modelle.}
  Große vortrainierte Modelle (ImageNet-CNNs, DINOv2, CLIP, MAE oder UniGaze)
  wurden für die Blickschätzung bisher vor allem mit vollständigem Fine-Tuning
  auf großen Datensätzen untersucht. Es fehlen Untersuchungen, wie gut sie sich
  mit nur wenigen Minuten Webcam-Kalibrierung an eine Person anpassen lassen.

  \item \textbf{Sparsame Anpassung für die Personalisierung.}
  Mehrere Arbeiten zeigen, dass wenige angepasste Parameter den größten Teil
  des persönlichen Fehlers erfassen. Trotzdem wurden reines Head-Training,
  Adapter und LoRA für die Personalisierung in der Blickschätzung kaum
  systematisch verglichen.

  \item \textbf{Labelqualität kurzer Kalibrierungen.}
  Reaktionszeit und Zielwechsel erzeugen Kalibrierbilder mit falschem Label.
  Bisherige Arbeiten behandeln falsche Labels meist nachher in großen
  Datensätzen oder verwerfen feste Zeitfenster. Wie sich ein genaues,
  protokollbasiertes Entfernen solcher Bilder auf die Personalisierung
  auswirkt, wurde noch nicht untersucht.

  \item \textbf{Realistische Evaluation unter Webcam-Bedingungen.}
  Viele Ergebnisse stammen aus Benchmarks mit kontrollierten Bedingungen oder
  aus zufälligen Splits auf Bildebene, die die Generalisierung überschätzen
  können. Wie Kamera, Brille, Hintergrund und Kalibrierdauer die
  Personalisierung beeinflussen, ist kaum dokumentiert.
\end{enumerate}

\subsection{Was diese Arbeit beiträgt}
\label{subsec:ueberleitung}
Diese Arbeit bearbeitet die Lücken mit einer einfachen und reproduzierbaren
Personalisierungspipeline für Webcam-Eye-Tracking:

\begin{itemize}
  \item \textbf{Eingefrorener Backbone mit kleinem Regression Head (F1, F2):}
        Ein auf ImageNet vortrainiertes ResNet-18 bleibt eingefroren. Seine
        Features werden gespeichert, und nur ein kleiner Regression Head wird pro
        Person trainiert. Ich vergleiche außerdem, ob das Trainieren weiterer
        Schichten hilft. Andere Backbones (DINOv2, CLIP, MAE, UniGaze) und
        Adapter oder LoRA sind der nächste Schritt (Kapitel~6).
  \item \textbf{Phasenprotokoll und Bereinigung (F3):} Das Aufnahmeprogramm
        speichert die Phasen \code{WAIT} und \code{MOVE}. So lässt sich die
        Wirkung des Entfernens von Transition Frames kontrolliert messen:
        Modelle mit ursprünglichen und mit bereinigten Daten werden unter sonst
        gleichen Bedingungen verglichen.
  \item \textbf{Systematische Ablationen:} Batch Size, Learning Rate, trainierbare
        Schichten, Activation, Loss Function (auch ein nach Position
        gewichteter L1-Loss) und Datenmenge werden mit wiederholten Läufen
        und Confidence Intervals untersucht.
  \item \textbf{Realistische Evaluation (F4):} Die Basismodelle werden mit dem
        personendisjunkten Split \code{norm\_subject} bewertet. Die
        Personalisierung wird auf 18 Sitzungen mit wechselnder Webcam (zwei Geräte),
        Brille, Hintergrund und Kalibrierdauer getestet, und zwar mit
        zufälligem und mit zeitlichem Split.
\end{itemize}

Kapitel~3 beschreibt die Methodik, Kapitel~4 die Ergebnisse. Kapitel~5
beantwortet damit die Forschungsfragen aus Kapitel~1.

% ==============================================================================
